
\documentclass[11pt,a4paper]{article}

\usepackage{fontspec}
\usepackage[french]{babel}
\usepackage{amsmath,amssymb,amsthm,mathtools}
\usepackage{geometry}
\usepackage{xcolor}
\usepackage{fancyhdr}
\usepackage{enumitem}
\usepackage{hyperref}
\usepackage{titlesec}
\usepackage{microtype}
\usepackage{booktabs}
\usepackage{array}
\usepackage{tabularx}
\usepackage{longtable}
\usepackage{colortbl}

\geometry{margin=2.2cm}
\setmainfont{DejaVu Serif}
\setsansfont{DejaVu Sans}
\setmonofont{DejaVu Sans Mono}

\definecolor{accent}{HTML}{1F4E79}
\definecolor{lightaccent}{HTML}{EAF2F8}
\definecolor{lightgray}{HTML}{F7F7F7}

\pagestyle{fancy}
\fancyhf{}
\lhead{\small Cours Deep Reinforcement Learning}
\chead{\small Classe : 4IA}
\rhead{\small Année : 2026}
\lfoot{\small Enseignant : Mourad Zéraï, PhD - ESPRIT School of Engineering}
\cfoot{}
\rfoot{\thepage}
\setlength{\headheight}{14pt}

\titleformat{\section}{\large\bfseries\color{accent}}{\thesection.}{0.5em}{}
\titleformat{\subsection}{\normalsize\bfseries\color{accent}}{\thesubsection.}{0.5em}{}
\titleformat{\subsubsection}{\normalsize\bfseries}{\thesubsubsection.}{0.5em}{}

\newtheorem*{definition}{Définition}
\newtheorem*{rappel}{Rappel}
\newtheorem*{hypothese}{Hypothèse}
\newtheorem*{propriete}{Propriété}
\newtheorem*{theoreme}{Théorème}
\newtheorem*{remarque}{Remarque}

\newcommand{\E}{\mathbb{E}}
\newcommand{\Pbb}{\mathbb{P}}
\newcommand{\R}{\mathbb{R}}
\newcommand{\argmax}{\operatorname*{argmax}}
\newcommand{\given}{\,|\,}

\begin{document}

\begin{center}
    {\LARGE\bfseries Dérivation et explication pas à pas de TD(0)}\\[0.4em]
    {\large Note de cours rigoureuse, détaillée et appliquée à FrozenLake}
\end{center}

\vspace{0.6em}

\section{Objectif}

Le but de cette note est d'expliquer \emph{TD(0)} pour l'évaluation d'une politique, de dériver sa règle de mise à jour pas à pas et de l'appliquer à FrozenLake.

\section{Cadre}

On fixe une politique $\pi$. On cherche à estimer la fonction de valeur
\[
v_{\pi}(s)=\E_{\pi}[G_t \given S_t=s].
\]

\begin{rappel}
L'équation de Bellman pour $v_{\pi}$ est
\[
v_{\pi}(s)=\sum_a \pi(a\given s)\sum_{s',r} p(s',r\given s,a)\left[r+\gamma v_{\pi}(s')\right].
\]
\end{rappel}

On veut maintenant estimer cette fonction sans connaître explicitement le modèle complet et sans attendre la fin des épisodes comme en Monte Carlo.

\section{Point de départ : la décomposition du retour}

On sait que
\[
G_t = R_{t+1}+\gamma G_{t+1}.
\]
En prenant l'espérance conditionnelle sous $\pi$, on obtient
\[
v_{\pi}(S_t)=\E_{\pi}[R_{t+1}+\gamma G_{t+1}\given S_t].
\]

\subsection{Approximation fondamentale}

L'idée de TD(0) est la suivante :
\begin{itemize}[leftmargin=1.6em]
    \item on ne remplace pas $G_{t+1}$ par le retour complet réellement observé ;
    \item on le remplace par l'estimation courante
    \[
    V(S_{t+1}).
    \]
\end{itemize}

On obtient alors une cible approchée :
\[
R_{t+1}+\gamma V(S_{t+1}).
\]

\begin{remarque}
Cette opération s'appelle souvent \emph{bootstrapping}. On utilise une estimation pour en améliorer une autre.
\end{remarque}

\section{Erreur temporelle}

On définit l'erreur TD
\[
\delta_t = R_{t+1}+\gamma V(S_{t+1})-V(S_t).
\]

\paragraph{Interprétation.}
\begin{itemize}[leftmargin=1.6em]
    \item Si $\delta_t>0$, la transition observée suggère que $V(S_t)$ est trop petite.
    \item Si $\delta_t<0$, elle suggère que $V(S_t)$ est trop grande.
\end{itemize}

\section{Mise à jour TD(0)}

Comme pour tout schéma de correction incrémentale, on déplace l'estimation actuelle dans la direction de la cible :
\[
V(S_t)\leftarrow V(S_t)+\alpha\delta_t.
\]
En remplaçant $\delta_t$, on obtient la règle classique :
\[
V(S_t)\leftarrow V(S_t)+\alpha\left[R_{t+1}+\gamma V(S_{t+1})-V(S_t)\right].
\]

\section{Justification étape par étape}

\subsection{Étape 1 : cible idéale}

Si l'on connaissait exactement $v_{\pi}$, alors pour toute transition moyenne sous la politique $\pi$,
\[
v_{\pi}(S_t)\approx R_{t+1}+\gamma v_{\pi}(S_{t+1})
\]
au sens de l'espérance conditionnelle.

\subsection{Étape 2 : remplacement de l'inconnu par l'estimation courante}

Comme $v_{\pi}$ est inconnue, on remplace $v_{\pi}(S_{t+1})$ par $V(S_{t+1})$.

\subsection{Étape 3 : correction proportionnelle à l'erreur}

On applique ensuite le schéma général
\[
x_{\text{new}} = x_{\text{old}}+\alpha(\text{cible}-x_{\text{old}}),
\]
avec
\[
x_{\text{old}}=V(S_t), \qquad \text{cible}=R_{t+1}+\gamma V(S_{t+1}).
\]

\subsection{Étape 4 : obtention de TD(0)}

On obtient
\[
V(S_t)\leftarrow V(S_t)+\alpha\left[R_{t+1}+\gamma V(S_{t+1})-V(S_t)\right].
\]

\section{Pourquoi on parle de TD(0)}

\begin{itemize}[leftmargin=1.6em]
    \item \textbf{TD} signifie \emph{Temporal Difference}, car on corrige la valeur d'un état en comparant deux prédictions faites à des instants successifs.
    \item Le \textbf{0} signifie qu'on ne prend qu'un pas de bootstrap avant de revenir à l'estimation. On ne déroule pas plusieurs étapes avant la correction.
\end{itemize}

\section{Comparaison avec Monte Carlo}

\begin{itemize}[leftmargin=1.6em]
    \item Monte Carlo attend la fin de l'épisode et utilise le retour complet
    \[
    G_t.
    \]
    \item TD(0) met à jour immédiatement après chaque transition avec
    \[
    R_{t+1}+\gamma V(S_{t+1}).
    \]
\end{itemize}

\paragraph{Conséquence.}
TD(0) apprend plus vite en ligne, mais introduit un biais de bootstrap. Monte Carlo est sans bootstrap mais attend la fin de l'épisode.

\section{Convergence : idée générale}

Dans le cadre tabulaire, si :
\begin{itemize}[leftmargin=1.6em]
    \item la politique $\pi$ est fixée ;
    \item tous les états pertinents sont visités ;
    \item les pas $(\alpha_t)$ satisfont les conditions standard d'approximation stochastique ;
\end{itemize}
alors TD(0) converge vers $v_{\pi}$.

\paragraph{Lecture intuitive.}
La mise à jour est une approximation stochastique du point fixe de l'opérateur de Bellman associé à la politique $\pi$.

\section{Application à FrozenLake}

\subsection{Choix d'une politique}

On fixe par exemple :
\begin{itemize}[leftmargin=1.6em]
    \item une politique aléatoire uniforme ;
    \item une politique manuelle ;
    \item une politique greedy issue d'un autre algorithme.
\end{itemize}

\subsection{Déroulement}

À chaque transition observée dans FrozenLake :
\begin{enumerate}[leftmargin=1.8em]
    \item on part d'un état $S_t$ ;
    \item on choisit une action selon $\pi$ ;
    \item on observe $R_{t+1}$ et $S_{t+1}$ ;
    \item on met à jour
    \[
    V(S_t)\leftarrow V(S_t)+\alpha\left[R_{t+1}+\gamma V(S_{t+1})-V(S_t)\right].
    \]
\end{enumerate}

\subsection{Cas glissant}

Quand \texttt{is\_slippery=True}, les transitions sont stochastiques. TD(0) est alors particulièrement naturel, car il apprend directement à partir des transitions observées sans exiger le modèle complet.

\subsection{Lecture visuelle}

Dans une interface FrozenLake, on peut afficher :
\begin{itemize}[leftmargin=1.6em]
    \item la valeur estimée $V(s)$ dans chaque case ;
    \item la TD error de la dernière transition ;
    \item l'évolution de $V$ au fil des épisodes.
\end{itemize}

\section{Résumé}

TD(0) est l'algorithme le plus simple pour estimer $v_{\pi}$ en ligne. Il remplace le retour complet par une cible de bootstrap à un pas :
\[
R_{t+1}+\gamma V(S_{t+1}),
\]
puis corrige
\[
V(S_t)\leftarrow V(S_t)+\alpha\left[R_{t+1}+\gamma V(S_{t+1})-V(S_t)\right].
\]
Dans FrozenLake, il permet d'illustrer clairement l'idée de Temporal Difference, surtout si l'on compare sa vitesse de mise à jour à celle de Monte Carlo.

\end{document}
